\documentclass[10pt,a4paper]{amsart}
\usepackage{amsmath,amssymb,amsthm,graphicx}
\usepackage{verbatim}
\newenvironment{reference}{\comment}{\endcomment}
%\newenvironment{reference}{}{}
\newenvironment{slogan}{\comment}{\endcomment}
\newenvironment{history}{\comment}{\endcomment}

% For commutative diagrams we use Xy-pic
\usepackage[all]{xy}

% We use 2cell for 2-commutative diagrams.
\xyoption{2cell}
% Theorem environments.
%
\theoremstyle{plain}
\newtheorem{theorem}[subsection]{Theorem}
\newtheorem{proposition}[subsection]{Proposition}
\newtheorem{lemma}[subsection]{Lemma}

\theoremstyle{definition}
\newtheorem{definition}[subsection]{Definition}
\newtheorem{example}[subsection]{Example}
\newtheorem{exercise}[subsection]{Exercise}
\newtheorem{situation}[subsection]{Situation}

\theoremstyle{remark}
\newtheorem{remark}[subsection]{Remark}
\newtheorem{remarks}[subsection]{Remarks}

\numberwithin{equation}{subsection}

% Macros
%
\def\lim{\mathop{\mathrm{lim}}\nolimits}
\def\colim{\mathop{\mathrm{colim}}\nolimits}
\def\Spec{\mathop{\mathrm{Spec}}}
\def\Hom{\mathop{\mathrm{Hom}}\nolimits}
\def\Ext{\mathop{\mathrm{Ext}}\nolimits}
\def\SheafHom{\mathop{\mathcal{H}\!\mathit{om}}\nolimits}
\def\SheafExt{\mathop{\mathcal{E}\!\mathit{xt}}\nolimits}
\def\Sch{\mathit{Sch}}
\def\Mor{\mathop{\mathrm{Mor}}\nolimits}
\def\Ob{\mathop{\mathrm{Ob}}\nolimits}
\def\Sh{\mathop{\mathit{Sh}}\nolimits}
\def\NL{\mathop{N\!L}\nolimits}
\def\CH{\mathop{\mathrm{CH}}\nolimits}
\def\proetale{{pro\text{-}\acute{e}tale}}
\def\etale{{\acute{e}tale}}
\def\QCoh{\mathit{QCoh}}
\def\Ker{\mathop{\mathrm{Ker}}}
\def\Im{\mathop{\mathrm{Im}}}
\def\Coker{\mathop{\mathrm{Coker}}}
\def\Coim{\mathop{\mathrm{Coim}}}

% Boxtimes
%
\DeclareMathSymbol{\boxtimes}{\mathbin}{AMSa}{"02}

%
% Macros for moduli stacks/spaces
%
\def\QCohstack{\mathcal{QC}\!\mathit{oh}}
\def\Cohstack{\mathcal{C}\!\mathit{oh}}
\def\Spacesstack{\mathcal{S}\!\mathit{paces}}
\def\Quotfunctor{\mathrm{Quot}}
\def\Hilbfunctor{\mathrm{Hilb}}
\def\Curvesstack{\mathcal{C}\!\mathit{urves}}
\def\Polarizedstack{\mathcal{P}\!\mathit{olarized}}
\def\Complexesstack{\mathcal{C}\!\mathit{omplexes}}
% \Pic is the operator that assigns to X its picard group, usage \Pic(X)
% \Picardstack_{X/B} denotes the Picard stack of X over B
% \Picardfunctor_{X/B} denotes the Picard functor of X over B
\def\Pic{\mathop{\mathrm{Pic}}\nolimits}
\def\Picardstack{\mathcal{P}\!\mathit{ic}}
\def\Picardfunctor{\mathrm{Pic}}
\def\Deformationcategory{\mathcal{D}\!\mathit{ef}}
\usepackage{cleveref}
\crefname{originalreference}{original reference}{original references}
\Crefname{originalreference}{Original reference}{Original references}
\makeatletter
\newlabel{situation-versal}{{0DQU}{source}{}{sourcepointer-1}{}}
\newlabel{situation-versal@cref}{{[originalreference][0][]0DQU}{[1][0][]source}{}{}{}}
\makeatother
\crefname{originalreference}{original reference}{original references}
\Crefname{originalreference}{Original reference}{Original references}
\makeatletter
\newlabel{artin-lemma-effective}{{07X8}{source}{}{sourcepointer-2}{}}
\newlabel{artin-lemma-effective@cref}{{[originalreference][0][]07X8}{[1][0][]source}{}{}{}}
\makeatother
\crefname{originalreference}{original reference}{original references}
\Crefname{originalreference}{Original reference}{Original references}
\makeatletter
\newlabel{stacks-morphisms-lemma-noetherian-finite-type-finite-presentation}{{0DQJ}{source}{}{sourcepointer-3}{}}
\newlabel{stacks-morphisms-lemma-noetherian-finite-type-finite-presentation@cref}{{[originalreference][0][]0DQJ}{[1][0][]source}{}{}{}}
\makeatother
\crefname{originalreference}{original reference}{original references}
\Crefname{originalreference}{Original reference}{Original references}
\makeatletter
\newlabel{stacks-limits-proposition-characterize-locally-finite-presentation}{{0CMY}{source}{}{sourcepointer-4}{}}
\newlabel{stacks-limits-proposition-characterize-locally-finite-presentation@cref}{{[originalreference][0][]0CMY}{[1][0][]source}{}{}{}}
\makeatother
\crefname{originalreference}{original reference}{original references}
\Crefname{originalreference}{Original reference}{Original references}
\makeatletter
\newlabel{artin-lemma-approximate-versal}{{07XH}{source}{}{sourcepointer-5}{}}
\newlabel{artin-lemma-approximate-versal@cref}{{[originalreference][0][]07XH}{[1][0][]source}{}{}{}}
\makeatother
\crefname{originallemma}{original lemma}{original lemmas}
\Crefname{originallemma}{Original lemma}{Original lemmas}
\makeatletter
\newlabel{lemma-compare-versal-ring-completion}{{0DQZ}{source}{}{sourcepointer-6}{}}
\newlabel{lemma-compare-versal-ring-completion@cref}{{[originallemma][0][]0DQZ}{[1][0][]source}{}{}{}}
\makeatother
\title{Artin approximation realizes a versal ring as a completed local ring of a smooth morphism}
\author{The Stacks Project Authors}
\date{stacks-project:stacks-geometry.tex:local-snapshot-20261012; GNU FDL 1.2 or later}
\begin{document}
\maketitle
\noindent\textit{Editorial note.} All delivered mathematical text is the author's own, copied verbatim from the pinned Stacks Project chapter file stacks-geometry.tex; no statement or proof step was rewritten, reordered or removed.\par\medskip
\noindent\textit{Editorial note.} Span recomputation by label pairing (LANE\_BRIEF 10.2): in this chapter the author's proof environment follows the statement environment as a standalone block; the spec covers the complete environment pair (statement span, proof span) as recorded in the pool candidate.\par\medskip
\noindent\textit{Editorial note.} Occupation scan: combined claim\_key + statement\_anchor + generic-line scan over all registration files (active and withdrawn) found 0 hits for stacks-geometry.tex (ge-occupied.json, 2026-10-14); 207 route=geometry / chapter-list field hits in focused-seeds are alias noise (route name and alphabetical chapter enumerations), excluded per the ladder alias-blind-spot caveat.\par\medskip
\ifdefined\hypertarget\hypertarget{sourcepointer-1}{}\fi
\noindent\textit{Source reference.} Original reference 0DQU: sources/stacks-chow/tags.txt, source lines 17199--17199. Cross-chapter standard prerequisite cited by the author inside the counted span; it is used as a pointer and its own proof is not part of this item.\par\medskip
\ifdefined\hypertarget\hypertarget{sourcepointer-2}{}\fi
\noindent\textit{Source reference.} Original reference 07X8: sources/stacks-chow/tags.txt, source lines 10068--10068. Cross-chapter standard prerequisite cited by the author inside the counted span; it is used as a pointer and its own proof is not part of this item.\par\medskip
\ifdefined\hypertarget\hypertarget{sourcepointer-3}{}\fi
\noindent\textit{Source reference.} Original reference 0DQJ: sources/stacks-chow/tags.txt, source lines 17189--17189. Cross-chapter standard prerequisite cited by the author inside the counted span; it is used as a pointer and its own proof is not part of this item.\par\medskip
\ifdefined\hypertarget\hypertarget{sourcepointer-4}{}\fi
\noindent\textit{Source reference.} Original reference 0CMY: sources/stacks-chow/tags.txt, source lines 15873--15873. Cross-chapter standard prerequisite cited by the author inside the counted span; it is used as a pointer and its own proof is not part of this item.\par\medskip
\ifdefined\hypertarget\hypertarget{sourcepointer-5}{}\fi
\noindent\textit{Source reference.} Original reference 07XH: sources/stacks-chow/tags.txt, source lines 10077--10077. Cross-chapter standard prerequisite cited by the author inside the counted span; it is used as a pointer and its own proof is not part of this item.\par\medskip
\ifdefined\hypertarget\hypertarget{sourcepointer-6}{}\fi
\noindent\textit{Source reference.} Original lemma 0DQZ: sources/stacks-chow/tags.txt, source lines 17204--17204. Cross-chapter standard prerequisite cited by the author inside the counted span; it is used as a pointer and its own proof is not part of this item.\par\medskip
\begin{lemma}
\label{lemma-Artin-approximation-by-smooth-morphism}
In Situation \ref{situation-versal}. Let $x_0 : \Spec(k) \to \mathcal{X}$
be a morphism such that $\Spec(k) \to S$ is of finite type with image $s$.
Let $A$ be a versal ring to $\mathcal{X}$ at $x_0$.
If $\mathcal{O}_{S, s}$ is a G-ring, then we may find a smooth morphism
$U \to \mathcal{X}$ whose source is a scheme and a point
$u_0 \in U$ with residue field $k$, such that
\begin{enumerate}
\item $\Spec(k) \to U \to \mathcal{X}$ coincides with the given morphism $x_0$,
\item there is an isomorphism $\mathcal{O}_{U, u_0}^\wedge \cong A$.
\end{enumerate}
\end{lemma}

\begin{proof}
Let $(\xi_n, f_n)$ be the versal formal object over $A$.
By Artin's Axioms, Lemma \ref{artin-lemma-effective}
we know that $\xi = (A, \xi_n, f_n)$ is effective.
By assumption $\mathcal{X}$ is locally of finite presentation over $S$
(use Morphisms of Stacks, Lemma
\ref{stacks-morphisms-lemma-noetherian-finite-type-finite-presentation}),
and hence limit preserving by Limits of Stacks, Proposition
\ref{stacks-limits-proposition-characterize-locally-finite-presentation}.
Thus Artin approximation as in
Artin's Axioms, Lemma \ref{artin-lemma-approximate-versal}
shows that we may find a morphism $U \to \mathcal{X}$ with
source a finite type $S$-scheme, containing a point $u_0 \in U$
of residue field $k$ satisfying (1) and (2) such that $U \to \mathcal{X}$
is versal at $u_0$. By Lemma \ref{lemma-compare-versal-ring-completion}
after shrinking $U$ we may assume $U \to \mathcal{X}$ is smooth.
\end{proof}

\end{document}
